\documentclass[12pt]{article}
\usepackage{amsmath}
\usepackage{amssymb}
\usepackage{graphicx}
\usepackage{caption}
\usepackage{float}
\usepackage[utf8]{inputenc}
\usepackage[T1]{fontenc}
\usepackage{geometry}
\usepackage{parskip}
\usepackage[table]{xcolor}
\usepackage{tikz}
\usetikzlibrary{positioning,arrows.meta,shapes.geometric,fit}
\usepackage[english]{babel}
\usepackage{tabularx}
\usepackage{array}
\usepackage{colortbl}
\usepackage{enumitem}
\usepackage[hidelinks]{hyperref}
\usepackage{indentfirst}
\geometry{a4paper, margin=1in}

\usepackage{listings}
\lstset{
    basicstyle=\ttfamily\footnotesize,
    breaklines=true,
    frame=single,
    backgroundcolor=\color{gray!10},
    keywordstyle=\color{blue!70!black},
    commentstyle=\color{green!50!black},
    stringstyle=\color{red!60!black},
    showstringspaces=false,
    language=C
}

\definecolor{stellantis}{RGB}{26,58,92}
\definecolor{sectionblue}{RGB}{44,95,138}
\definecolor{ddbrownlight}{RGB}{255,243,230}
\definecolor{ddbrown}{RGB}{180,100,30}

\usepackage[most]{tcolorbox}
\tcbuselibrary{breakable}

\renewcommand{\arraystretch}{1.3}

\begin{document}

\begin{center}
{\LARGE\bfseries\color{stellantis} Relatório Consolidado de V\&V}\\[6pt]
{\large Verificação integral do módulo \textbf{CAN Bus} para ASIL-D}\\[10pt]
{\small
\begin{tabular}{rl}
\textbf{Projeto:}           & Autonomous Emergency Braking (AEB) --- UFPE / Stellantis \\
\textbf{Validador (cross):} & Lourenço \\
\textbf{Autor original:}    & Renato Fagundes \\
\textbf{Escopo:}            & \texttt{aeb\_can.c}, \texttt{aeb\_can.h} (Communication / CAN Bus) \\
\textbf{Padrão aplicado:}   & ISO 26262-6:2018 · ASIL-D · MISRA C:2012 · CAN 2.0A · DBC \\
\textbf{Data:}              & 23 de abril de 2026 \\
\end{tabular}
}
\end{center}

\vspace{0.3cm}

\begin{tcolorbox}[colback=ddbrownlight, colframe=ddbrown, title=\textbf{Sumário executivo}]
Este documento consolida as cinco atividades de verificação aplicadas ao módulo CAN Bus durante a
cross-validation independente para ASIL-D: testes unitários nominais, cobertura estrutural MC/DC, injeção
sistemática de falhas, verificação dinâmica de segurança de memória, e análise estática. O processo
atingiu \textbf{100\% em todas as métricas estruturais} (statement, branch, MC/DC) com a combinação da
suíte nominal existente (autoria Renato) e duas suítes complementares produzidas durante a
cross-validation: testes estruturais de gaps de cobertura e testes de injeção de falhas. A verificação
dinâmica (Valgrind, ASan) não detectou erros de segurança de memória. A injeção sistemática de falhas
(31~asserções em 4~categorias) \textbf{teve todas as asserções aprovadas}; contudo, o UBSan
(\texttt{-fsanitize=float-cast-overflow}) revelou \textbf{uma vulnerabilidade semântica} na função
auxiliar \texttt{encode\_unsigned}, catalogada neste relatório com o respectivo patch proposto.
\end{tcolorbox}

% ═══════════════════════════════════════════════════════════════════════════
\section{Visão geral das atividades}

A ISO~26262-6:2018 define, nas tabelas~8, 10, 11 e 12 da Parte~6, um conjunto de técnicas de verificação
``altamente recomendadas'' (++) para ASIL-D. Aplicamos cinco dessas técnicas ao escopo CAN:

\begin{table}[H]
\centering
\caption{Panorama das atividades de V\&V para o módulo CAN.}
\footnotesize
\setlength{\tabcolsep}{4pt}
\begin{tabularx}{\textwidth}{|c|>{\raggedright\arraybackslash}p{3.1cm}
                              |>{\raggedright\arraybackslash}p{3.1cm}
                              |>{\raggedright\arraybackslash}p{2.7cm}
                              |>{\raggedright\arraybackslash}X|}
\hline
\rowcolor{sectionblue}
\textcolor{white}{\textbf{\#}} &
\textcolor{white}{\textbf{Atividade}} &
\textcolor{white}{\textbf{Ferramenta}} &
\textcolor{white}{\textbf{Norma ISO 26262-6}} &
\textcolor{white}{\textbf{Resultado (escopo CAN)}} \\ \hline

1 & Testes unitários nominais & Framework próprio (C99) & Tab.~10 item~1b (++) & 13 testes · 36 asserções · 100\% PASS \\ \hline
2 & Cobertura MC/DC          & gcc-14 + gcov-14        & Tab.~12 item~1c (++) & 100\% stmt/branch/MC/DC \\ \hline
3 & Fault injection          & Framework próprio       & Tab.~11 item~1e (++) & 31 asserções · 31 PASS \\ \hline
4 & Memory safety dinâmica   & Valgrind + ASan + UBSan & Tab.~8 item~1d (++)  & Valgrind/ASan: 0; UBSan: 1 UB \\ \hline
5 & Análise estática         & cppcheck + MISRA addon  & Tab.~8 item~1b (++)  & 2 advisory em \texttt{aeb\_can.\{c,h\}} \\ \hline

\end{tabularx}
\end{table}

Todas as ferramentas são open-source, reproduzíveis em Ubuntu~24.04, e automatizadas pelo novo alvo
\texttt{make vv-can} adicionado ao Makefile raiz.

\subsection{Princípio de independência}

A ISO~26262 exige, para ASIL-D, que a verificação de unidade seja conduzida por engenheiro
independente do implementador (\S5.4.8). A equipe adota a regra de \textit{cross-assignment}: cada
membro valida um módulo que não implementou. O módulo CAN, originalmente implementado por
\textbf{Renato Fagundes}, foi validado independentemente por \textbf{Lourenço}. Este relatório documenta
integralmente essa validação.

\subsection{Delimitação de escopo}

\begin{tcolorbox}[colback=ddbrownlight, colframe=ddbrown,
                  title=\textbf{Escopo deste relatório vs. MISRA consolidada}]
A Seção~\ref{sec:misra} apresenta \textit{análise estática escopada em \texttt{aeb\_can.c} / \texttt{aeb\_can.h}}.
O \textbf{relatório MISRA consolidado do sistema} (cobrindo todos os módulos, requisito do Checklist~\#8
da Estratégia de Entrega Final) é um entregável separado, produzido pela branch de autoria do Renato
(mesma pessoa que é autor original do código CAN aqui validado --- os dois papéis são distintos e
não criam conflito de independência, pois este relatório valida \textit{apenas} \texttt{aeb\_can.\{c,h\}}).
Onde esses dois relatórios se sobrepuserem, a fonte-de-verdade para o escopo CAN é este relatório;
para o escopo sistêmico, o relatório consolidado.
\end{tcolorbox}

\subsection{Estrutura da suíte de testes}

Diferente do UDS (onde a suíte nominal do autor original atingiu 100\% de cobertura estrutural por
construção), o módulo CAN teve \textbf{gaps de cobertura} em ramos defensivos (switch statements) que
nenhum teste baseado em requisitos exercitava. Seguindo o mesmo padrão documentado para o módulo
PID/Alert, foi acrescentada uma suíte \textit{complementar estrutural} (\texttt{test\_can\_struct.c},
autoria Lourenço) com o mínimo de testes necessários para fechar os gaps. A estrutura final é:

\begin{table}[H]
\centering
\caption{Suítes de teste aplicadas ao módulo CAN.}
\footnotesize
\setlength{\tabcolsep}{4pt}
\begin{tabularx}{\textwidth}{|>{\raggedright\arraybackslash}p{4.2cm}
                              |>{\raggedright\arraybackslash}p{2.2cm}
                              |>{\centering\arraybackslash}p{1.6cm}
                              |>{\raggedright\arraybackslash}X|}
\hline
\rowcolor{sectionblue}
\textcolor{white}{\textbf{Suíte}} &
\textcolor{white}{\textbf{Autoria}} &
\textcolor{white}{\textbf{Casos}} &
\textcolor{white}{\textbf{Propósito}} \\ \hline
\texttt{test\_can.c}         & Renato   & 13 & Testes baseados em requisitos (FR-CAN-001..004) \\ \hline
\texttt{test\_can\_struct.c} & Lourenço & 13 & Fechamento estrutural de ramos defensivos \\ \hline
\texttt{test\_can\_fault.c}  & Lourenço & 16 & Injeção sistemática de falhas \\ \hline
\end{tabularx}
\end{table}

% ═══════════════════════════════════════════════════════════════════════════
\section{Atividade 1 --- Testes unitários nominais}

A suíte nominal, escrita pelo autor original, cobre os requisitos funcionais \textbf{FR-CAN-001} a
\textbf{FR-CAN-004}. São testes de aceitação por requisito.

\begin{table}[H]
\centering
\caption{Distribuição dos testes unitários nominais (escopo CAN, arquivo \texttt{test\_can.c}).}
\footnotesize
\begin{tabularx}{\textwidth}{|>{\raggedright\arraybackslash}p{5.5cm}
                              |>{\centering\arraybackslash}p{1.5cm}
                              |>{\raggedright\arraybackslash}X|}
\hline
\rowcolor{sectionblue}
\textcolor{white}{\textbf{Grupo de testes}} &
\textcolor{white}{\textbf{Casos}} &
\textcolor{white}{\textbf{Requisitos cobertos}} \\ \hline
Pack/unpack de sinais (round-trip) & 1 & FR-CAN-003 (DBC signal encoding) \\ \hline
Inicialização (\texttt{can\_init})  & 2 & FR-CAN-004 (baud rate 500~kbit/s, falha de HAL) \\ \hline
Decodificação RX AEB\_EgoVehicle (0x100)   & 2 & FR-CAN-002, FR-CAN-003 \\ \hline
Decodificação RX AEB\_RadarTarget (0x120)  & 1 & FR-CAN-002, FR-CAN-003 \\ \hline
Decodificação RX AEB\_DriverInput (0x101)  & 1 & FR-CAN-002, FR-CAN-003 \\ \hline
Timeout de RX (3~$\times$~20~ms)           & 2 & FR-CAN-002 (acceptance: 60~ms) \\ \hline
Transmissão AEB\_BrakeCmd (0x080)          & 1 & FR-CAN-003 \\ \hline
Transmissão AEB\_FSMState (0x200) a 50~ms  & 1 & FR-CAN-001 \\ \hline
Transmissão AEB\_Alert (0x300)             & 1 & FR-CAN-003 \\ \hline
Tratamento de falha de TX (HAL)            & 1 & NFR-SAF-ROB (robustez TX) \\ \hline
\textbf{Total}                             & \textbf{13} & \textbf{FR-CAN-001..004, NFR-SAF-ROB} \\ \hline
\end{tabularx}
\end{table}

\begin{tcolorbox}[colback=green!5, colframe=green!50!black]
\textbf{Resultado global:} 13 de 13 casos de teste passaram (\textbf{36 asserções,
100\% PASS}). Execução via \texttt{make test}, integrada à pipeline de CI do
GitHub~Actions (\texttt{ci.yml}).
\end{tcolorbox}

% ═══════════════════════════════════════════════════════════════════════════
\section{Atividade 2 --- Cobertura estrutural MC/DC}

\textbf{Modified Condition/Decision Coverage (MC/DC)} é a métrica de cobertura estrutural ``altamente
recomendada'' (++) pela ISO~26262-6 Tab.~12 para ASIL-D. Para cada condição atômica, MC/DC exige
demonstrar que, mantendo as demais fixas, alterar apenas ela muda o resultado da decisão --- prova de
independência lógica.

\subsection{Ferramenta utilizada}

Utilizou-se o suporte nativo do GCC~14 via flag \texttt{-fcondition-coverage} (introduzida em 2024),
combinado com \texttt{gcov-14 -{}-conditions}. Os dados de cobertura dos três binários de teste
(\texttt{test\_can}, \texttt{test\_can\_struct}, \texttt{test\_can\_fault}) são consolidados via
\texttt{gcov-tool-14 merge}. Um alvo Makefile \texttt{mcdc-can} automatiza o fluxo e grava artefatos em
\texttt{reports/vv\_can/coverage\_mcdc/}.

\subsection{Análise de gaps e testes complementares}

A suíte nominal isolada atinge 89{,}82\% de cobertura de statement para \texttt{aeb\_can.c}. Os 17
statements não executados concentram-se em dois \texttt{switch}~statements defensivos
(\texttt{FSM\_state}~$\rightarrow$~\texttt{brake\_mode} e
\texttt{FSM\_state}~$\rightarrow$~\texttt{TTC\_threshold}) e em ramos defensivos de \texttt{NULL}-guard
e \texttt{CAN\_ERR\_TX}. Dois esforços complementares, produzidos pelo validador cross (Lourenço),
fecharam os gaps:

\begin{itemize}[leftmargin=1.2em, itemsep=2pt]
\item \texttt{test\_can\_fault.c} cobriu os ramos de \texttt{NULL}-guard e de saturação do
      \texttt{alive\_counter}.
\item \texttt{test\_can\_struct.c} cobriu cada \texttt{case} dos switch defensivos e o ramo \texttt{CAN\_ERR\_TX} de \texttt{can\_tx\_fsm\_state}.
\end{itemize}

\subsection{Resultado final}

\begin{table}[H]
\centering
\caption{Cobertura estrutural do módulo CAN (linha de base + complementares, via \texttt{gcov-tool merge}).}
\footnotesize
\begin{tabular}{|l|c|c|l|}
\hline
\rowcolor{sectionblue}
\textcolor{white}{\textbf{Métrica}} &
\textcolor{white}{\textbf{Valor}} &
\textcolor{white}{\textbf{Meta ASIL-D}} &
\textcolor{white}{\textbf{Status}} \\ \hline
Statement coverage (\texttt{aeb\_can.c})  & 100{,}00\% (167/167) & $\geq$95\% & $\checkmark$ CONFORME \\ \hline
Branch coverage (\texttt{aeb\_can.c})     & 100{,}00\% (81/81)   & $\geq$95\% & $\checkmark$ CONFORME \\ \hline
MC/DC coverage (condition outcomes)       & 100{,}00\% (68/68)   & $\geq$95\% & $\checkmark$ CONFORME \\ \hline
Calls exercised                           & 100{,}00\% (41/41)   & ---        & --- \\ \hline
\end{tabular}
\end{table}

\begin{tcolorbox}[colback=ddbrownlight, colframe=ddbrown,
                  title=\textbf{Observação de implementação}]
O módulo CAN contém duas áreas de alta complexidade ciclomática: o \texttt{switch} sobre
\texttt{fsm\_state} em \texttt{can\_tx\_brake\_cmd} (7~\texttt{case}~+ \texttt{default}) e em
\texttt{can\_tx\_fsm\_state} (3~\texttt{case}~+ \texttt{default}). A suíte nominal do autor só
exercitava dois desses cases; a suíte estrutural complementar é o que torna a cobertura 100\%.
Isto mantém a \textbf{intencionalidade do autor original} preservada (a suíte dele continua sendo a
\textit{fonte de verdade} para requisitos) e adiciona, por fora, a varredura estrutural que a ISO~26262
exige para ASIL-D.
\end{tcolorbox}

\textbf{Reprodução:} \texttt{make mcdc-can} gera \texttt{run\_\{nominal,fault,struct\}.log},
\texttt{aeb\_can.c.gcov} e relatório HTML em \texttt{reports/vv\_can/coverage\_mcdc/}.

% ═══════════════════════════════════════════════════════════════════════════
\section{Atividade 3 --- Fault injection sistemática}
\label{sec:fault}

Fault injection alimenta o código com entradas inválidas, corrompidas ou extremas para verificar que o
comportamento permanece fail-safe. É ``altamente recomendada'' (++) pela ISO~26262-6 Tab.~11 item~1e
para ASIL-D.

\subsection{Categorias exercitadas}

\begin{table}[H]
\centering
\caption{Categorias de falha aplicadas ao módulo CAN.}
\footnotesize
\setlength{\tabcolsep}{4pt}
\begin{tabularx}{\textwidth}{|c|>{\raggedright\arraybackslash}p{3.5cm}
                              |>{\raggedright\arraybackslash}X|}
\hline
\rowcolor{sectionblue}
\textcolor{white}{\textbf{Cat.}} &
\textcolor{white}{\textbf{Natureza}} &
\textcolor{white}{\textbf{Exemplos aplicados}} \\ \hline
A & Valores inválidos  & Ponteiros \texttt{NULL} em todas as APIs públicas; DLC inferior ao esperado por ID; ID desconhecido; falha forçada do HAL de init (15 asserções) \\ \hline
B & Extremos numéricos & \texttt{NaN}, $\pm\infty$, \texttt{FLT\_MAX}, overflow do cast em \texttt{encode\_unsigned}; \texttt{raw\_value} excedendo a largura do sinal em \texttt{can\_pack\_signal} (7 asserções) \\ \hline
C & Estado corrompido (SEU) & \texttt{alive\_counter}~$=$~0xAA; \texttt{rx\_miss\_count}~$=$~0xFF; \texttt{tx\_cycle\_counter}~$=$~\texttt{UINT32\_MAX}; \texttt{initialised}~$=$~0x55 (5 asserções) \\ \hline
D & Persistência       & Ciclo timeout $\rightarrow$ RX válido $\rightarrow$ timeout; 100~toggles de \texttt{force\_send\_fail}; burst TX~>~32~frames (4 asserções) \\ \hline
E & Temporização       & \textbf{Fora de escopo} no nível unitário (bancada host sem scheduler); diferido para SIL \\ \hline
\end{tabularx}
\end{table}

\subsection{Predicado de fail-safe}

A predicate de ``fail-safe'' usada nas categorias A e B é:
\emph{sob entrada anormal, a função afetada deve (i) retornar um código de erro explícito
(\texttt{CAN\_ERR\_TX} ou \texttt{CAN\_ERR\_INIT}), ou (ii) produzir um frame CAN em que o valor
bruto empacotado esteja contido dentro da largura declarada do sinal, sem corromper sinais
adjacentes.} Qualquer comportamento que dereferencie \texttt{NULL}, escreva bits fora da janela do
sinal, ou transmita bytes implementation-defined \textbf{não é fail-safe}.

\subsection{Resultado}

\begin{table}[H]
\centering
\caption{Resultado fault injection, módulo CAN.}
\scriptsize
\setlength{\tabcolsep}{4pt}
\begin{tabular}{|l|c|c|c|c|l|}
\hline
\rowcolor{sectionblue}
\textcolor{white}{\textbf{Categoria}} &
\textcolor{white}{\textbf{Testes}} &
\textcolor{white}{\textbf{Asserções}} &
\textcolor{white}{\textbf{Passaram}} &
\textcolor{white}{\textbf{Falharam}} &
\textcolor{white}{\textbf{Veredito}} \\ \hline
A -- Valores inválidos  & 2  & 15 & 15 & 0 & Robusto \\ \hline
B -- Extremos numéricos & 1  & 7  & 7  & 0 & Por coincidência (ver nota) \\ \hline
C -- SEU / estado       & 1  & 5  & 5  & 0 & Robusto \\ \hline
D -- Persistência       & 1  & 4  & 4  & 0 & Robusto \\ \hline
\textbf{TOTAL}          & \textbf{5} & \textbf{31} & \textbf{31} & \textbf{0} & --- \\ \hline
\end{tabular}
\end{table}

\begin{tcolorbox}[colback=red!5, colframe=red!50!black,
                  title=\textbf{Nota sobre a Categoria~B (extremos numéricos)}]
As sete asserções de \texttt{NaN}/$\pm\infty$/\texttt{FLT\_MAX} passaram o predicado de fail-safe, mas
\textit{por coincidência do comportamento do GCC~14 em x86-64}: o cast \texttt{(uint32\_t)NaN} e
\texttt{(uint32\_t)}$\pm\infty$ produz 0 nessa toolchain/plataforma. O padrão C11 \S6.3.1.4 classifica
esses casts como comportamento indefinido, e o UBSan (com \texttt{-fsanitize=float-cast-overflow})
confirma isso na linha 89 (Seção~\ref{sec:memory}). \textbf{A robustez observada não é garantida} ao
mudar compilador, target ou nível de otimização. O patch do Bug~\#1 (Seção~\ref{sec:bug1}) torna o
comportamento fail-safe por projeto, eliminando a dependência do acaso.
\end{tcolorbox}

\textbf{Reprodução:} \texttt{make fault-can} grava \texttt{reports/vv\_can/fault\_injection/run.log}.

% ═══════════════════════════════════════════════════════════════════════════
\section{Atividade 4 --- Memory safety}
\label{sec:memory}

Verificação dinâmica de segurança de memória usando três ferramentas complementares, aplicadas às
suítes nominal e de fault injection do módulo CAN.

\begin{table}[H]
\centering
\caption{Ferramentas de verificação dinâmica.}
\footnotesize
\begin{tabularx}{\textwidth}{|>{\raggedright\arraybackslash}p{2.6cm}
                              |>{\raggedright\arraybackslash}X
                              |>{\raggedright\arraybackslash}p{3.0cm}|}
\hline
\rowcolor{sectionblue}
\textcolor{white}{\textbf{Ferramenta}} &
\textcolor{white}{\textbf{Classes detectadas}} &
\textcolor{white}{\textbf{Mecanismo}} \\ \hline
Valgrind 3.22    & Memory leaks, uso de memória não inicializada, acesso fora de limites de blocos alocados & Interpretação em VM \\ \hline
AddressSanitizer & Heap/stack/global buffer overflow, use-after-free, use-after-scope, double-free         & Instrumentação em compile-time \\ \hline
UBSanitizer      & Signed integer overflow, division by zero, invalid shift, misaligned pointer, \textbf{float-cast-overflow} & Instrumentação em compile-time \\ \hline
\end{tabularx}
\end{table}

\subsection{Resultado}

\begin{table}[H]
\centering
\caption{Resultado das verificações dinâmicas (suítes CAN).}
\footnotesize
\begin{tabular}{|l|l|l|l|}
\hline
\rowcolor{sectionblue}
\textcolor{white}{\textbf{Binário}} &
\textcolor{white}{\textbf{Valgrind}} &
\textcolor{white}{\textbf{ASan}} &
\textcolor{white}{\textbf{UBSan (float-cast-overflow)}} \\ \hline
test\_can (nominal)       & 0 erros & 0 erros & 0 erros \\ \hline
test\_can\_fault          & 0 erros & 0 erros & \textbf{1 local de UB detectado} \\ \hline
\end{tabular}
\end{table}

\begin{lstlisting}[language={}, caption={Diagnóstico UBSan encontrado na suíte de fault injection.}]
src/communication/aeb_can.c:89:32: runtime error: nan is outside the range of
    representable values of type 'unsigned int'
\end{lstlisting}

UBSan confirma independentemente o que a Categoria~B do fault injection já sugeria por análise: o cast
\texttt{(uint32\_t)(raw\_f + 0.5F)} na função auxiliar \texttt{encode\_unsigned} é comportamento
indefinido quando \texttt{raw\_f} é \texttt{NaN}. As chamadas afetadas são todas as encodings DBC de
campos float na transmissão (pressão de freio, ângulo TTC, acelerações).

Valgrind e ASan limpos: o código não acessa memória indevidamente, não vaza, e não usa memória
não-inicializada. A vulnerabilidade encontrada é \textbf{semântica} (cast UB), não de memória.

\textbf{Reprodução:} \texttt{make memory-can} grava \texttt{valgrind\_\{nominal,fault\}.log} e
\texttt{ubsan\_\{nominal,fault\}.log} em \texttt{reports/vv\_can/memory\_safety/}.

% ═══════════════════════════════════════════════════════════════════════════
\section{Atividade 5 --- Análise estática}
\label{sec:misra}

A análise estática é exigida pela Tab.~8 item~1b da ISO~26262-6 (uso de subconjunto da linguagem ---
MISRA~C:2012 é o padrão de-facto no setor automotivo).

\subsection{Ferramenta e escopo}

Utilizou-se o \texttt{cppcheck 2.13} com o addon MISRA, via alvo Makefile \texttt{misra-can}. O escopo é
estritamente \texttt{src/communication/aeb\_can.c} e \texttt{include/aeb\_can.h}. A análise também
inclui os cabeçalhos transitivamente incluídos (\texttt{aeb\_types.h}, \texttt{aeb\_config.h}) mas os
achados nesses arquivos pertencem ao escopo da \textit{análise MISRA consolidada do sistema}
(não este relatório; ver Seção~1.2).

\subsection{Resultado}

\begin{table}[H]
\centering
\caption{Resultado cppcheck+MISRA (apenas escopo CAN, \texttt{aeb\_can.\{c,h\}}).}
\footnotesize
\begin{tabular}{|l|c|l|}
\hline
\rowcolor{sectionblue}
\textcolor{white}{\textbf{Severidade}} &
\textcolor{white}{\textbf{Issues}} &
\textcolor{white}{\textbf{Comentário}} \\ \hline
error       & 0 & Nenhum erro semântico \\ \hline
warning     & 0 & Nenhum comportamento suspeito \\ \hline
performance & 0 & Nenhum antipadrão \\ \hline
portability & 0 & Nenhum uso não-portável \\ \hline
style       & 2 & Advisory --- ver detalhe abaixo \\ \hline
\textbf{Total escopo CAN} & \textbf{2} & \textbf{Apenas advisories, sem high/medium} \\ \hline
\end{tabular}
\end{table}

\subsection{Detalhe das advisories in-scope}

\begin{itemize}[leftmargin=1.2em, itemsep=4pt]
\item \texttt{aeb\_can.c:313} --- MISRA C:2012 \textbf{Rule 10.6} (advisory, style):
      \texttt{uint32\_t brake\_req = (pid\_out->brake\_pct > 0.0F) ? 1U : 0U;}
      A expressão ternária é classificada pela regra como uma composite expression atribuída a tipo
      mais largo que o tipo essencial. \textbf{Decisão de justificar} em vez de modificar: o
      comportamento está correto e a mudança de forma (e.g. via variável intermediária tipada)
      agregaria verbosidade sem benefício funcional.
\item \texttt{aeb\_can.h:45} --- MISRA C:2012 \textbf{Rule 2.5} (advisory, style):
      macro \texttt{CAN\_OK} (e parceiras \texttt{CAN\_ERR\_*}) marcada como possivelmente não usada
      dentro do escopo do header. Na prática são constantes públicas usadas por múltiplos módulos e
      pelos testes, mas o cppcheck não consegue verificar isso sem a análise global.
\end{itemize}

Para transparência de processo, a varredura completa também visita os cabeçalhos transitivamente
incluídos. Foram reportados \textbf{60 diagnósticos \textit{style}} nesses cabeçalhos compartilhados
(\texttt{aeb\_types.h} e \texttt{aeb\_config.h}), todos do grupo MISRA 2.3/2.4/2.5 --- regras
\textit{advisory} sobre declarações não-utilizadas, atribuíveis à estrutura superset dos cabeçalhos
do sistema. Serão tratados no relatório consolidado MISRA sistêmico.

\begin{tcolorbox}[colback=green!5, colframe=green!50!black]
\textbf{Conclusão:} o módulo CAN não contém erros, warnings, performance ou portability issues MISRA
dentro do seu escopo de código próprio. As duas advisories de style têm justificativa documentada e
não impedem a certificação. Zero violações high/medium.
\end{tcolorbox}

% ═══════════════════════════════════════════════════════════════════════════
\section{Bug encontrado e patch proposto}
\label{sec:bug1}

\subsection{Bug \#1 --- Cast float$\to$inteiro sem validação em \texttt{encode\_unsigned}}

\textbf{Localização:} \texttt{src/communication/aeb\_can.c:89}

\textbf{Descrição:} a função auxiliar \texttt{encode\_unsigned} converte um valor físico
(\texttt{float32\_t}) em valor bruto (\texttt{uint32\_t}) para empacotamento DBC. O código atual:

\begin{lstlisting}[caption={Código atual (vulnerável).}]
static uint32_t encode_unsigned(float32_t physical,
                                float32_t factor,
                                float32_t offset)
{
    float32_t raw_f = (physical - offset) / factor;
    uint32_t  raw   = 0U;

    if (raw_f < 0.0F) {
        raw = 0U;
    } else {
        raw = (uint32_t)(raw_f + 0.5F);  /* <-- UB se raw_f = NaN ou Inf */
    }
    return raw;
}
\end{lstlisting}

\textbf{Problema:} se \texttt{physical} for \texttt{NaN}, a comparação \texttt{raw\_f < 0.0F} é sempre
\texttt{false} (regra IEEE~754), portanto o fluxo vai para o \texttt{else} e executa
\texttt{(uint32\_t)NaN}, que é \textit{undefined behaviour} segundo C11~\S6.3.1.4. O mesmo vale para
\texttt{physical = +$\infty$}: \texttt{raw\_f} é \texttt{+$\infty$}, \texttt{raw\_f + 0.5F} é
\texttt{+$\infty$}, e o cast é UB. Para \texttt{physical = }\texttt{FLT\_MAX}, o valor de \texttt{raw\_f}
pode exceder \texttt{UINT32\_MAX} e o cast é UB.

\textbf{Por que ainda funciona em produção:} GCC~14 em x86-64 materializa todos esses UB como
\texttt{0}, que por acaso é um valor \textit{fail-safe} (freio zero, TTC zero). \textbf{Isto não é
garantido} em outra toolchain, nível de otimização, ou processador alvo (ARM Cortex).

\textbf{Impacto:} em compilação diferente, o frame CAN AEB\_BrakeCmd ou AEB\_FSMState pode ser
transmitido com um valor bruto arbitrário (lixo de registro FPU). O ECU de freio a jusante
interpretaria isso como pressão hidráulica aleatória --- potencial frenagem não comandada.

\subsubsection*{Patch proposto}

\begin{lstlisting}[caption={Normalização explícita de NaN/Inf e saturação.}]
#include <math.h>
#include <stdint.h>

static uint32_t encode_unsigned(float32_t physical,
                                float32_t factor,
                                float32_t offset)
{
    uint32_t raw = 0U;

    /* Guard NaN and +/-Inf per C11 6.3.1.4 */
    if (isnan(physical) || isinf(physical)) {
        raw = 0U;
    } else {
        float32_t raw_f = (physical - offset) / factor;
        if (raw_f < 0.0F) {
            raw = 0U;
        } else if (raw_f > (float32_t)UINT32_MAX) {
            raw = UINT32_MAX;                /* Saturate to max uint */
        } else {
            raw = (uint32_t)(raw_f + 0.5F);  /* safe path */
        }
    }
    return raw;
}
\end{lstlisting}

\textbf{Teste de aceitação:} o binário \texttt{test\_can\_fault\_san} (ASan+UBSan) deve executar com
\texttt{0~runtime error}. Os 31~asserções da suíte de fault injection devem continuar passando, e os
13~testes nominais do Renato devem continuar passando sem modificação.

% ═══════════════════════════════════════════════════════════════════════════
\section{Plano de ação e cronograma}

\begin{table}[H]
\centering
\caption{Ações necessárias para fechar a V\&V do CAN.}
\footnotesize
\setlength{\tabcolsep}{4pt}
\begin{tabularx}{\textwidth}{|c|>{\raggedright\arraybackslash}X
                              |>{\raggedright\arraybackslash}p{2.8cm}
                              |>{\centering\arraybackslash}p{1.6cm}|}
\hline
\rowcolor{sectionblue}
\textcolor{white}{\textbf{\#}} &
\textcolor{white}{\textbf{Ação}} &
\textcolor{white}{\textbf{Responsável}} &
\textcolor{white}{\textbf{Esforço}} \\ \hline
1 & Aplicar patch do Bug~\#1 em \texttt{aeb\_can.c:89} (\texttt{encode\_unsigned}) & Renato   & 1--2 h \\ \hline
2 & Re-executar \texttt{make vv-can} e confirmar UBSan 0 runtime error             & Lourenço & 10 min \\ \hline
3 & Confirmar que 13 nominais + 13 estruturais + 31 fault continuam 100\% PASS     & Lourenço & 5 min  \\ \hline
4 & Abrir PR com o patch em \texttt{fix/can-encode-robustness}, base \texttt{development} & Renato & 30 min \\ \hline
5 & Preencher linhas FR-CAN-* do Excel Funcional                                   & Lourenço & 1 h    \\ \hline
6 & Abrir issue sistêmica para validação NaN/Inf em outros encoders float (PID, TTC) & Lourenço / Renato & 30 min \\ \hline
\end{tabularx}
\end{table}

\subsection{Narrativa para a defesa}

\begin{tcolorbox}[colback=gray!5, colframe=gray!50]
\itshape
``Aplicando a stack completa de verificação para ASIL-D ao módulo CAN, medimos cobertura estrutural
100\% (statement, branch, MC/DC) com o auxílio de duas suítes complementares --- estrutural e de
fault injection --- à suíte nominal do autor original. Valgrind e AddressSanitizer não encontraram
erros de memória. A injeção sistemática de falhas --- 31 asserções em 4 categorias --- passou todas
as asserções sob a toolchain alvo, mas a categoria de extremos numéricos só é fail-safe por
coincidência. UBSan (\texttt{-fsanitize=float-cast-overflow}) confirmou independentemente uma
vulnerabilidade de cast float$\to$inteiro em \texttt{encode\_unsigned}. O patch proposto transforma
comportamento fail-safe \textit{por coincidência da toolchain} em fail-safe \textit{por projeto}. O
ciclo completo detecção $\to$ correção $\to$ re-validação independente exemplifica a aplicação
prática do princípio de independência da ISO~26262-6 \S5.4.8.''
\end{tcolorbox}

% ═══════════════════════════════════════════════════════════════════════════
\section{Rastreabilidade: requisito $\to$ teste $\to$ evidência}

\begin{table}[H]
\centering
\caption{Sumário de rastreabilidade por requisito. Detalhe completo (uma linha por teste) em \texttt{functional\_tests\_can.csv} --- fonte para a aba CAN do Excel Funcional.}
\footnotesize
\setlength{\tabcolsep}{4pt}
\begin{tabularx}{\textwidth}{|>{\raggedright\arraybackslash}p{2.3cm}
                              |>{\centering\arraybackslash}p{1.4cm}
                              |>{\raggedright\arraybackslash}X
                              |>{\centering\arraybackslash}p{1.3cm}|}
\hline
\rowcolor{sectionblue}
\textcolor{white}{\textbf{Requisito}} &
\textcolor{white}{\textbf{\# Testes}} &
\textcolor{white}{\textbf{Evidência principal}} &
\textcolor{white}{\textbf{P/F}} \\ \hline
FR-CAN-001 Ego dynamics TX (50 ms) & 2  & \texttt{test\_can.test\_tx\_fsm\_period} + struct HAL-fail    & P \\ \hline
FR-CAN-002 Radar data RX + timeout & 4  & \texttt{test\_rx\_radar\_target}, \texttt{test\_rx\_timeout*} & P \\ \hline
FR-CAN-003 DBC signal encoding     & 6  & \texttt{test\_signal\_roundtrip}, \texttt{test\_tx\_brake\_cmd}, \texttt{test\_tx\_alert}, \texttt{test\_rx\_ego\_vehicle}, \texttt{test\_rx\_driver\_input}, \texttt{test\_rx\_yaw\_rate} & P \\ \hline
FR-CAN-004 Baud rate 500 kbit/s    & 2  & \texttt{test\_init\_success}, \texttt{test\_init\_failure}    & P \\ \hline
NFR-SAF-ROB Robustez               & 31 & \texttt{fault\_injection/run.log} + UBSan                     & \textbf{F} (Bug~\#1) \\ \hline
\end{tabularx}
\end{table}

% ═══════════════════════════════════════════════════════════════════════════
\section{Conformidade com ISO 26262 ASIL-D}

\begin{table}[H]
\centering
\caption{Mapeamento explícito das atividades $\to$ tabelas normativas ISO 26262-6:2018.}
\footnotesize
\setlength{\tabcolsep}{4pt}
\begin{tabular}{|c|c|l|c|>{\raggedright\arraybackslash}p{6cm}|}
\hline
\rowcolor{sectionblue}
\textcolor{white}{\textbf{Tabela}} &
\textcolor{white}{\textbf{Item}} &
\textcolor{white}{\textbf{Técnica exigida}} &
\textcolor{white}{\textbf{ASIL-D}} &
\textcolor{white}{\textbf{Aplicada por}} \\ \hline
8  & 1b & Uso de subconjunto da linguagem    & ++ & MISRA C:2012 via cppcheck (Seção~\ref{sec:misra}) \\ \hline
8  & 1d & Implementação defensiva            & ++ & Valgrind + ASan + UBSan (Seção~\ref{sec:memory}) \\ \hline
8  & 1e & Uso restrito de memória dinâmica   & ++ & Valgrind: 0 bytes heap \\ \hline
10 & 1b & Testes baseados em requisitos      & ++ & \texttt{test\_can.c} (13 testes FR-CAN-*) \\ \hline
11 & 1e & Injeção de falhas / error guessing & ++ & \texttt{test\_can\_fault.c} (31 asserções, 4 categorias) \\ \hline
12 & 1a & Statement coverage                 & ++ & gcov-14 $\to$ 100\% \\ \hline
12 & 1b & Branch coverage                    & ++ & gcov-14 $\to$ 100\% \\ \hline
12 & 1c & MC/DC coverage                     & ++ & gcc-14 \texttt{-fcondition-coverage} $\to$ 100\% \\ \hline
\end{tabular}
\end{table}

\begin{tcolorbox}[colback=green!5, colframe=green!50!black]
\textbf{Conclusão de conformidade:} todas as 8~técnicas estruturais marcadas como ``altamente
recomendadas'' (++) para ASIL-D nas tabelas~8, 10, 11 e 12 da ISO~26262-6 foram aplicadas ao módulo
CAN, com evidência documentada e reproduzível via \texttt{make vv-can}. A única não-conformidade
encontrada (Bug~\#1) está rastreada, com patch proposto e teste de aceitação definido.
\end{tcolorbox}

% ═══════════════════════════════════════════════════════════════════════════
\section*{Referências}

\begin{itemize}[leftmargin=1.2em, itemsep=2pt]
\item ISO 26262-6:2018 --- Road vehicles --- Functional safety --- Part 6: Product development at the software
      level. \S5.4.8 (independence), Tables 8, 10, 11, 12.
\item ISO 11898-1:2015 --- Road vehicles --- Controller area network (CAN) --- Part 1: Data link layer and
      physical signalling.
\item MISRA C:2012, 3rd edition (2019). Rules 2.3/2.4/2.5, 10.6 (advisory; shared-header and ternary scope).
\item IEEE 754-2019 --- Standard for Floating-Point Arithmetic. Comportamento de NaN em operações
      aritméticas e comparações.
\item ISO/IEC 9899:2011 (C11) \S6.3.1.4 --- Real floating to integer conversions.
\item GCC 14 Release Notes --- Condition coverage support (\texttt{-fcondition-coverage}).
\item Valgrind Developers. Valgrind User Manual, version 3.22.
\item Clang/LLVM Team. \texttt{-fsanitize=float-cast-overflow} documentation.
\item Ammann, P. \& Offutt, J. (2016). Introduction to Software Testing, 2nd ed.
\item AEB-Stellantis Project --- \texttt{final\_delivery\_strategy.tex}, abril 2026.
\item Fagundes, R. --- \texttt{Relatorio\_Consolidado\_VV\_UDS.tex}, abril 2026 (template de referência).
\end{itemize}

\vspace{0.5cm}
\noindent
{\footnotesize\itshape
Relatório produzido como parte da cross-validation independente da Deliverable~4/4 do projeto AEB ---
Residência Tecnológica em Desenvolvimento de Software Embarcado para o Setor Automotivo · UFPE /
Stellantis · 2026. Artefatos reproduzíveis via \texttt{make vv-can}. Commit:
\texttt{\$\{GIT\_SHA\}} (preencher no tag de release).
}

\end{document}
